\begin{recipe}{Broccoli uit de oven}
\kicker[Bijgerecht,Vegetarisch,Doordeweeks]{Bijgerecht · vegetarisch · doordeweeks}
\index[register]{Broccoli uit de oven}
\index[register]{Broccoli}
\index[register]{Kaas}

\meta{25 minuten \dvd Oven 200\,°C}

\lettrine{B}{roccoli} uit de oven is een simpel bijgerecht: kort voorgekookte
roosjes met een laag geraspte kaas, een paar minuten de oven in tot de kaas
gesmolten is. Handig naast aardappelen, een andere groente of een stuk vlees,
en zo op tafel op een doordeweekse avond.

\blockrule{Ingrediënten}
\begin{ingredients}
  \ing{1 kg}{broccoli}\index[register]{Broccoli}
  \ing{150 g}{jong belegen geraspte kaas}\index[register]{Kaas}
  \ingb{versgemalen peper}
\end{ingredients}

\blockrule{Bereiding}
\tip{Laat de broccoli na het afgieten goed uitstomen voor je hem in de
ovenschaal legt, bijvoorbeeld even terug in de droge pan. Blijft er te veel
vocht in de broccoli achter, dan wordt de ovenschaal nat en verdunt dat de
gesmolten kaas.}
\begin{steps}
  \item Verwarm de oven voor op 200\,°C (boven- en onderwarmte).
  \item Was de broccoli en verdeel in roosjes.
  \item Kook de broccoliroosjes 5 tot 6 minuten in ruim water met zout tot
        ze beetgaar zijn.
  \item Giet de broccoli af en laat goed uitstomen (zie tip).
  \item Verdeel de broccoli over een ovenschaal (ca.\ 26\,×\,18\,×\,7 cm),
        strooi de geraspte kaas erover en bestrooi met versgemalen peper.
  \item Bak 10 tot 15 minuten in de oven tot de kaas gesmolten is en
        lichtbruine plekken heeft.
\end{steps}
\end{recipe}
